\documentclass[../../main.tex]{subfiles}
\begin{document}
\section{积分曲线}\label{section:ch09-积分曲线}

假设 $M$ 是带边或不带边的光滑流形。若 $\gamma\colon J\to M$ 是光滑曲线，则对每个 $t\in J$，速度向量 $\gamma'(t)$ 是 $T_{\gamma(t)}M$ 中的向量。本节描述一种反向操作的方法：给定每一点处的切向量，我们寻找一条曲线，使其每一点处的速度等于该点处给定的向量。

若 $V$ 是 $M$ 上的向量场，则 $V$ 的一条\textbf{积分曲线}是一条可微曲线 $\gamma\colon J\to M$，使其每一点处的速度等于 $V$ 在该点处的值：
\begin{align*}
\gamma'(t)=V_{\gamma(t)}\quad\text{对所有 }t\in J.
\end{align*}
若 $0\in J$，则点 $\gamma(0)$ 称为 $\gamma$ 的\textbf{起点}。（"积分曲线"这一术语的理由稍后解释。注意这是一个需要某种可微性假设的定义，因为对于仅连续的曲线，积分曲线的定义将毫无意义。）

\begin{figure}[ht]
\centering
\includegraphics[width=0.8\textwidth]{fig-09-1.png}
\caption{向量场的积分曲线}\label{figure:ch09-integral-curve}
\end{figure}

\begin{example}\label{example:ch09-01}[积分曲线]
\begin{enumerate}[(a)]
\item 设 $(x,y)$ 是 $\mathbb{R}^2$ 上的标准坐标，$V=\partial/\partial x$ 是第一个坐标向量场。容易验证 $V$ 的积分曲线恰为平行于 $x$ 轴的直线，其参数化为 $\gamma(t)=(a+t,b)$，其中 $a,b$ 为常数。因此，平面上每一点都有唯一的积分曲线从该点出发，不同积分曲线的像或者相同，或者不相交。

\item 设 $W=x\partial/\partial y-y\partial/\partial x$ 是 $\mathbb{R}^2$ 上的向量场。若 $\gamma\colon\mathbb{R}\to\mathbb{R}^2$ 是光滑曲线，在标准坐标下写作 $\gamma(t)=(x(t),y(t))$，则 $\gamma$ 为积分曲线的条件 $\gamma'(t)=W_{\gamma(t)}$ 转化为
\begin{align*}
x'(t)\frac{\partial}{\partial x}\bigg|_{\gamma(t)}+y'(t)\frac{\partial}{\partial y}\bigg|_{\gamma(t)}
=x(t)\frac{\partial}{\partial y}\bigg|_{\gamma(t)}-y(t)\frac{\partial}{\partial x}\bigg|_{\gamma(t)}.
\end{align*}
比较这些向量的分量，可知这等价于如下常微分方程组：
\begin{align*}
x'(t)=-y(t),\qquad y'(t)=x(t).
\end{align*}
这些方程的解为
\begin{align*}
x(t)=a\cos t-b\sin t,\qquad y(t)=a\sin t+b\cos t,
\end{align*}
其中 $a,b$ 为任意常数，因此每条形如 $\gamma(t)=(a\cos t-b\sin t,a\sin t+b\cos t)$ 的曲线都是 $W$ 的积分曲线。当 $(a,b)=(0,0)$ 时，这是常值曲线 $\gamma(t)\equiv(0,0)$；否则，它是一条逆时针环绕的圆。由于 $\gamma(0)=(a,b)$，我们再次看到每一点 $(a,b)\in\mathbb{R}^2$ 都有唯一的积分曲线从该点出发，各积分曲线的像或者相同，或者不相交。
\end{enumerate}
\end{example}

\begin{figure}[ht]
\centering
\includegraphics[width=0.8\textwidth]{fig-09-2.png}
\caption{平面中的向量场及其积分曲线}\label{figure:ch09-integral-curves-plane}
\end{figure}

如上第二个例子所示，寻找积分曲线归结为在光滑坐标卡中求解常微分方程组。假设 $V$ 是 $M$ 上的光滑向量场，$\gamma\colon J\to M$ 是光滑曲线。在光滑坐标域 $U\subseteq M$ 上，可以在局部坐标中将 $\gamma$ 写作 $\gamma(t)=(\gamma^1(t),\dots,\gamma^n(t))$。则 $\gamma$ 为 $V$ 的积分曲线的条件 $\gamma'(t)=V_{\gamma(t)}$ 可写为
\begin{align*}
\dot{\gamma}^i(t)\frac{\partial}{\partial x^i}\bigg|_{\gamma(t)}=V^i\bigl(\gamma(t)\bigr)\frac{\partial}{\partial x^i}\bigg|_{\gamma(t)},
\end{align*}
这归结为如下自治常微分方程组（ODE）：
\begin{align}
\dot{\gamma}^1(t)&=V^1\bigl(\gamma^1(t),\dots,\gamma^n(t)\bigr),\nonumber\\
&\vdots\label{align:ch09-01}\\
\dot{\gamma}^n(t)&=V^n\bigl(\gamma^1(t),\dots,\gamma^n(t)\bigr).\nonumber
\end{align}
（当上标使撇号难以辨认时，我们用点表示关于 $t$ 的普通导数。）此类方程组的基本事实是存在性、唯一性和光滑性定理，即\cref{theorem:appD-01}。（这就是"积分曲线"这一术语的理由，因为求解常微分方程组通常被称为"积分"该方程组。）稍后将推导该定理的详细推论；现在先注意如下简单结果。

\begin{proposition}\label{proposition:ch09-02}
设 $V$ 是光滑流形 $M$ 上的光滑向量场。对每个点 $p\in M$，存在 $\varepsilon>0$ 和光滑曲线 $\gamma\colon(-\varepsilon,\varepsilon)\to M$，它是从 $p$ 出发的 $V$ 的积分曲线。
\end{proposition}

\begin{proof}
这正是\cref{theorem:appD-01}的存在性陈述应用于 $V$ 的坐标表示。
\end{proof}

下面两个引理说明仿射重参数化如何影响积分曲线。

\begin{lemma}\label{lemma:ch09-03}[重标度引理]
设 $V$ 是光滑流形 $M$ 上的光滑向量场，$J\subseteq\mathbb{R}$ 是区间，$\gamma\colon J\to M$ 是 $V$ 的积分曲线。对任意 $a\in\mathbb{R}$，由 $\widetilde{\gamma}(t)=\gamma(at)$ 定义的曲线 $\widetilde{\gamma}\colon\widetilde{J}\to M$ 是向量场 $aV$ 的积分曲线，其中 $\widetilde{J}=\{t:at\in J\}$。
\end{lemma}

\begin{proof}
一种看法是作为局部坐标中链式法则的直接应用。更不变地，我们可以考察 $\widetilde{\gamma}'(t)$ 对定义在某点 $\widetilde{\gamma}(t_0)$ 邻域内的光滑实值函数 $f$ 的作用。由链式法则以及 $\gamma$ 是 $V$ 的积分曲线这一事实，
\begin{align*}
\widetilde{\gamma}'(t_0)f=\left.\frac{\mathrm{d}}{\mathrm{d}t}\right|_{t=t_0}(f\circ\widetilde{\gamma})(t)=\left.\frac{\mathrm{d}}{\mathrm{d}t}\right|_{t=t_0}(f\circ\gamma)(at)=a(f\circ\gamma)'(at_0)=a\gamma'(at_0)f=aV_{\widetilde{\gamma}(t_0)}f.
\end{align*}
\end{proof}

\begin{lemma}\label{lemma:ch09-04}[平移引理]
设 $V,M,J,\gamma$ 如前一引理。对任意 $b\in\mathbb{R}$，由 $\widehat{\gamma}(t)=\gamma(t+b)$ 定义的曲线 $\widehat{\gamma}\colon\widehat{J}\to M$ 也是 $V$ 的积分曲线，其中 $\widehat{J}=\{t:t+b\in J\}$。
\end{lemma}

\begin{exercise}\label{exercise:ch09-05}
证明平移引理。
\end{exercise}

\begin{proposition}\label{proposition:ch09-06}[积分曲线的自然性]
假设 $M$ 和 $N$ 是光滑流形，$F\colon M\to N$ 是光滑映射。则 $X\in\mathfrak{X}(M)$ 和 $Y\in\mathfrak{X}(N)$ 是 $F$-相关的，当且仅当 $F$ 将 $X$ 的积分曲线映为 $Y$ 的积分曲线，即对 $X$ 的每条积分曲线 $\gamma$，$F\circ\gamma$ 是 $Y$ 的积分曲线。
\end{proposition}

\begin{figure}[ht]
\centering
\includegraphics[width=0.8\textwidth]{fig-09-3.png}
\caption{$F$-相关向量场的流}\label{figure:ch09-flows-f-related}
\end{figure}

\begin{proof}
首先假设 $X$ 和 $Y$ 是 $F$-相关的，$\gamma\colon J\to M$ 是 $X$ 的积分曲线。若定义 $\sigma\colon J\to N$ 为 $\sigma=F\circ\gamma$，则
\begin{align*}
\sigma'(t)=(F\circ\gamma)'(t)=\mathrm{d}F_{\gamma(t)}(\gamma'(t))=\mathrm{d}F_{\gamma(t)}(X_{\gamma(t)})=Y_{F(\gamma(t))}=Y_{\sigma(t)},
\end{align*}
故 $\sigma$ 是 $Y$ 的积分曲线。

反过来，假设 $F$ 将 $X$ 的积分曲线映为 $Y$ 的积分曲线。给定 $p\in M$，令 $\gamma\colon(-\varepsilon,\varepsilon)\to M$ 是从 $p$ 出发的 $X$ 的积分曲线。由于 $F\circ\gamma$ 是从 $F(p)$ 出发的 $Y$ 的积分曲线，有
\begin{align*}
Y_{F(p)}=(F\circ\gamma)'(0)=\mathrm{d}F_p(\gamma'(0))=\mathrm{d}F_p(X_p),
\end{align*}
这表明 $X$ 和 $Y$ 是 $F$-相关的。
\end{proof}

\end{document}